\documentclass{beamer}
\usetheme{Warsaw}
\usecolortheme{structure}
\usepackage{polski}
\usepackage[utf8]{inputenc}
\usepackage{hyperref}
\usepackage{tikz}

\usepackage{listings}
\lstloadlanguages{Ruby}
\lstset{basicstyle=\ttfamily\color{black}\scriptsize,
        commentstyle = \ttfamily\color{red},
        keywordstyle=\ttfamily\color{blue},
        stringstyle=\color{orange},
        showstringspaces=false,
        breaklines=true}

\hypersetup{colorlinks=true}

\title{Ruby on Rails}
\subtitle{z czym to się je?}
\author{Marek Waligórski}
\date{31 stycznia 2015}
\institute{InfoMEET Poznań}
\titlegraphic{\includegraphics[width=2cm,height=2cm]{Ruby_logo.png}}

% Dodaj obrazek do kazdego frametitle
\setbeamertemplate{frametitle}{%
    \nointerlineskip
    \begin{beamercolorbox}[sep=0.3cm,ht=1.8em,wd=\paperwidth]{frametitle}
        \vbox{}
        \vspace{-1.2ex}
        \strut\insertframetitle\strut
        \hfill
        \includegraphics[height=0.5cm]{Ruby_logo.png}
    \end{beamercolorbox}
}

\newcommand{\myfullimage}[1]{%
    \hspace*{-1.1cm}
    \includegraphics[width=\paperwidth]{#1}
}

\newcommand{\crossout}[0]{%
\begin{tikzpicture}[remember picture,overlay,decoration=penciline]
\draw[decorate,line width=15pt,red!60!black]
  (current page.north west) -- (current page.south east);
\draw[decorate,line width=15pt,red!60!black]
  (current page.north east) -- (current page.center) -- (current page.south west);
\end{tikzpicture}
}

\newcommand{\myurl}[1]{%
    \begin{block}{}
       \centering \huge
       \includegraphics[width=0.35cm]{link.png}
       \hspace{0.1cm}
       \url{#1}
    \end{block}
}
\newcommand{\mytitle}[1]{%
    \begin{frame}[plain]
        \begin{centering}
        \vspace{1em}\par
            \begin{beamercolorbox}[rounded=true,sep=4pt,center]{part title}
               \huge #1
            \end{beamercolorbox}
        \end{centering}
    \end{frame}
}
\newcommand{\mysection}[1]{%
    \section{#1}
    \mytitle{#1}
}

\begin{document}

\begin{frame}[plain]
\titlepage
\end{frame}

\begin{frame}[plain]
\frametitle{Ruby on Rails = $?$}
\begin{block}{Ruby}
    Język programowania
\end{block}
\begin{block}{Rails}
    Framework MVC
\end{block}
\begin{block}{Narzędzia}
    I kilka przydatnych bibliotek
\end{block}
\end{frame}

\mysection{Ruby}

\begin{frame}[fragile]\frametitle{Ruby}
    \begin{block}{Język programowania}
    \begin{itemize}
        \item skryptowy (interpretowany)
        \item zorientowany obiektowo
        \item dynamicznie typowany
    \end{itemize}
    \end{block}
\end{frame}

\mytitle{\centering \huge \color{orange} Übung macht den Meister \par}

\subsection{Jak zacząć?}
\begin{frame}[fragile]
\frametitle{Wypróbuj Ruby'ego bez instalacji}
    \myurl{tryruby.org}
    \myfullimage{tryruby.png}
\end{frame}
\begin{frame}[fragile]
\frametitle{Koduj albo giń!}
    \myurl{www.bloc.io/ruby-warrior}
    \myfullimage{rubywarrior.png}
\end{frame}

\mytitle{Dlaczego warto poznać Ruby'ego? \par \vspace{2cm}
\includegraphics[width=3cm]{Ruby_logo.png}
}

\mytitle{Elegancja}

\begin{frame}[fragile]
\frametitle{Elegancja}
\begin{block}{Ruby brzmi jak prawdziwy język}
\begin{lstlisting}[language=Ruby,basicstyle=\tiny\ttfamily]
exit unless "restaurant".include? "aura"

['toast', 'cheese', 'wine'].each { |food| print food.capitalize }

population = 12_000_000_000

(1..10).include? 5
\end{lstlisting}
\end{block}
\end{frame}

%%%% OBIEKTY %%%%%
\mytitle{Wszystko jest obiektem}

\begin{frame}[fragile]
\frametitle{Wszystko jest obiektem}
\begin{block}{}
\begin{lstlisting}[language=Ruby]
5.class  # => Fixnum
5 + 3    # => 8
5.+(3)   # => 8
\end{lstlisting}
\end{block}
\end{frame}

\mytitle{Nic nie jest wieczne}
\begin{frame}[fragile]
\frametitle{Nic nie jest wieczne}
\begin{block}{}
\begin{lstlisting}[language=Ruby]
5.class  # => Fixnum
5 + 3    # => 8
5.+(3)   # => 8
\end{lstlisting}
\end{block}
\begin{block}{Psikus!}
\begin{lstlisting}[language=Ruby]
class Fixnum
  def +(x)
    self - x
  end
end

5 + 3    # => 2
\end{lstlisting}
\end{block}
\end{frame}

%%%% STALE %%%%%

\mytitle{Język dynamiczny}
\mytitle{Bardzo dynamiczny}
\begin{frame}[fragile]
\frametitle{Nic nie jest wieczne}
\begin{block}{Stałe}
\begin{lstlisting}[language=Ruby]
CONSTANT = 5
\end{lstlisting}
\end{block}
\end{frame}

\begin{frame}[fragile]
\frametitle{Nic nie jest wieczne}
\begin{block}{Stałe}
\begin{lstlisting}[language=Ruby]
CONSTANT = 5
\end{lstlisting}
\end{block}
\begin{block}{Psikus!}
\begin{lstlisting}[language=Ruby]
CONSTANT = 5
CONSTANT = 0
# warning: already initialized constant CONSTANT
puts CONSTANT
# => 0
\end{lstlisting}
\end{block}
\end{frame}

\mytitle{With great power comes great responsibility.}

\subsection{Zalety}
\mytitle{Czy Ruby ma jakieś zalety?}
\mytitle{Ma.}

%%%% ZMIENNE %%%

\begin{frame}[fragile]
\frametitle{Elegancja}
\begin{block}{Pętla}
\begin{lstlisting}[language=C++]
for (int i=0; i<5; i++) {
  printf("Hello!\n");
}
\end{lstlisting}
\end{block}
\only<2->{\crossout}
\end{frame}

\mytitle{Nie twórzmy niepotrzebnych zmiennych}

\begin{frame}[fragile]
\frametitle{Elegancja}
\begin{columns}
 \column{5cm}
    \begin{block}{Pętla C++}
\begin{lstlisting}[language=C++]
for (int i=0; i<5; i++) {
  printf("Hello!\n");
}
\end{lstlisting}
    \end{block}
 \column{5cm}
    \begin{block}{Pętla w Ruby (bez zmiennej)}
\begin{lstlisting}[language=Ruby]
5.times { puts "Hello!" }
\end{lstlisting}
    \end{block}
\end{columns}
\end{frame}

\mytitle{Zmienna grzecznie czeka na drugim planie}

\begin{frame}[fragile]
\frametitle{Elegancja}
\begin{block}{Ruby nie narzuca ograniczeń}
\begin{lstlisting}[language=Ruby]
5.times { |i| puts i }

0
1
2
3
4
\end{lstlisting}
\end{block}
\end{frame}

\mytitle{Jedna linijka}

\begin{frame}[fragile]
\frametitle{Elegancja}
\begin{block}{Podnieś do kwadratu wszystkie liczby w tablicy}
\begin{lstlisting}[language=Ruby]
tablica = []
for i in [1,2,3,4]
  tablica << i ** 2
end
tablica   # => [1,4,9,16]
\end{lstlisting}
\end{block}
\only<2->{\crossout}
\end{frame}

\begin{frame}[fragile]
\frametitle{Elegancja}
\begin{block}{Podnieś do kwadratu wszystkie liczby w tablicy}
\begin{lstlisting}[language=Ruby]
(1..4).map { |i| i ** 2 }  # => [1,4,9,16]
\end{lstlisting}
\end{block}
\end{frame}

\begin{frame}[fragile]
\frametitle{Elegancja}
\begin{block}{Sortowanie według ostatniej litery imienia?}
\begin{lstlisting}[language=Ruby,basicstyle=\tiny\ttfamily]
["Stefan", "Tomek", "Agata", "Maja"].sort_by { |imie| imie[-1] }

# => ["Maja", "Agata", "Tomek", "Stefan"]
\end{lstlisting}
\end{block}
\end{frame}

%%%%% TURNIEJ %%%%%%
\begin{frame}[plain]
\myfullimage{foosball.jpg}
\end{frame}

\begin{frame}[fragile]
\frametitle{Elegancja}
\begin{block}{Losowanie drużyn do turnieju}
\begin{lstlisting}[language=Ruby]
gracze = ["Ela", "Ola", "Agnieszka",
          "Marek", "Wojtek", "Marcin"]
\end{lstlisting}
\end{block}
\end{frame}
\begin{frame}[fragile]
\frametitle{Elegancja}
\begin{block}{Losowanie drużyn do turnieju}
\begin{lstlisting}[language=Ruby]
gracze = ["Ela", "Ola", "Agnieszka",
          "Marek", "Wojtek", "Marcin"]

gracze.shuffle.each_slice(2).to_a

# => [["Agnieszka", "Marek"],
#     ["Wojtek", "Ola"],
#     ["Marcin", "Ela"]]
\end{lstlisting}
\end{block}
\end{frame}

%%%% LICZBY %%%%

\begin{frame}[fragile]
\frametitle{Duże liczby?}
\begin{block}{No problem!}
\begin{lstlisting}[language=Ruby]
7 ** 200
  => 1046183829131435717501889961181681
365981918855017023365995014008403512576
742426225177438261490936405029306524825
254631417406318034368359118815075426733
9816534637456120001

7.class
# => Fixnum
(7 ** 200).class
# => Bignum
\end{lstlisting}
\end{block}
\end{frame}


%%%% METAPROGRAMMING %%%%%%
\mytitle{Meta-programowanie}

\begin{frame}[fragile]
\frametitle{Metaprogramowanie}
\begin{block}{Sztuczne metody}
\begin{lstlisting}[language=Ruby]
class Auto
  def wypisz_predkosc
    puts "predkosc = #{predkosc}"
  end

  def wypisz_moc
    puts "moc = #{moc}"
  end

  private

  def predkosc
    "200 km/h"
  end

  def moc
    120
  end
end
\end{lstlisting}
\end{block}
\only<2->{\crossout}
\end{frame}
\begin{frame}[plain]
    \centering \includegraphics[height=\paperheight]{dry_yoda.jpg} \par
\end{frame}
\mytitle{obiekt.send(nazwa\_metody)}
\begin{frame}[fragile]
\frametitle{Metaprogramowanie}
\begin{block}{Wywołanie wybranej metody}
\begin{lstlisting}[language=Ruby]
class Auto
  def wypisz(co)
    puts "#{co} = #{send(co)}"
  end

  private

  def predkosc
    "200 km/h"
  end

  def moc
    120
  end
end

a = Auto.new
a.wypisz('moc')
# moc = 120

\end{lstlisting}
\end{block}
\end{frame}

\begin{frame}[plain]
\begin{block}{}
\begin{centering}
\color{violet}Ruby
   
\only<1->{\color{black}potrafi wywołać metodę}
   
\only<2->{która}
   
\only<3->{\huge \color{red} NIE ISTNIEJE}

\end{centering}
\end{block}
\end{frame}

\begin{frame}[fragile]
\frametitle{Metaprogramowanie}
\begin{block}{Wywołanie nieistniejącej metody}
\begin{lstlisting}[language=Ruby]
class A
  def hello
    puts "Hello, world!"
  end
end

a = A.new
a.hello
# "Hello, world!"
a.inna_metoda
# NoMethodError: undefined method `inna_metoda' for a:A
\end{lstlisting}
\end{block}
\end{frame}

\mytitle{Obiekty są samoświadome 
\vspace{1cm}
\includegraphics[width=7cm]{gerty.jpg}}

\begin{frame}[fragile]
\frametitle{Metaprogramowanie}
\begin{block}{Pomocny program}
\begin{lstlisting}[language=Ruby]
class A
  def hello
    puts "Hello, world!"
  end

  private

  def method_missing(name, *args)
    puts "Czy miales na mysli: #{self.methods.join(', ')}"
  end
end
\end{lstlisting}
\end{block}
\end{frame}
\begin{frame}[fragile]
\frametitle{Metaprogramowanie}
\begin{block}{Pomocny program}
\begin{lstlisting}[language=Ruby,basicstyle=\tiny\ttfamily]
class A
  def hello
    puts "Hello, world!"
  end

  private

  def method_missing(name, *args)
    puts "Czy miales na mysli: #{self.methods.join(', ')}"
  end
end

a = A.new
a.inna_metoda
# Czy miales na mysli: hello, nil?, ===, =~, !~, eql?, hash, <=>, class, singleton_class, clone, dup, taint, tainted?, untaint, untrust, untrusted?, trust, freeze, frozen?, to_s, inspect, methods, singleton_methods, protected_methods, private_methods, public_methods, instance_variables, instance_variable_get, instance_variable_set, instance_variable_defined?, remove_instance_variable, instance_of?, kind_of?, is_a?, tap, send, public_send, respond_to?, extend, display, method, public_method, singleton_method, define_singleton_method, object_id, to_enum, enum_for, ==, equal?, !, !=, instance_eval, instance_exec, __send__, __id__
\end{lstlisting}
\end{block}
\end{frame}
\begin{frame}[fragile]
\frametitle{Metaprogramowanie}
\begin{block}{Trochę bardziej pomocny program}
\begin{lstlisting}[language=Ruby,basicstyle=\tiny\ttfamily]
class A
  def hello
    puts "Hello, world!"
  end

  private

  def method_missing(name, *args)
    metody = self.methods - Object.methods
    puts "Czy miales na mysli: #{metody.join(', ')}"
  end
end

a = A.new
a.inna_metoda
# Czy miales na mysli: hello
\end{lstlisting}
\end{block}
\end{frame}

\mytitle{To nie wszystko!}
\begin{frame}[fragile]
\frametitle{Biblioteki}
    \myurl{rubygems.org}
    \myfullimage{rubygems.png}
\end{frame}

\mytitle{95,447 gemów \footnotesize(29/01/2015)}

\begin{frame}[fragile]
\frametitle{Biblioteki}
\begin{block}{Gemfile}
\begin{lstlisting}[language=Ruby,basicstyle=\tiny\ttfamily]
source 'https://rubygems.org'

gem 'rails', '3.2.21'
gem 'capistrano', '>= 2.15',require: false
gem 'netguru', :git => 'git://github.com/netguru/netguru.git'

gem 'rollbar', '1.4.0'
gem 'rails-backbone'
gem 'bson_ext'
gem 'carrierwave', '~> 0.9.0'
gem 'carrierwave-mongoid', require: 'carrierwave/mongoid'
gem 'colorize', require: false
gem 'dalli'
gem 'daemons'
gem 'decent_exposure'
gem 'draper'
...
\end{lstlisting}
\end{block}
\end{frame}

%\begin{frame}[fragile]
%\frametitle{implementacje}
    %* MRI
    %- JRuby - różnice
%\end{frame}


\begin{frame}[fragile]
\frametitle{Co nam dał?}

- Technologie sieciowe - where it excels
* lessCSS / SASS
* HAML
* CoffeeScript
* Rails!!!
\end{frame}

\subsection{Z czym się tego nie je}

\begin{frame}[fragile]
\frametitle{Z czym się tego nie je?}
    * GameDev
    * High-performance tasks (ale można integrować z modułami C++, ale nie ma np.\ bytecode'u - tylko AST)
    * Wątki - sztuczne (JRuby obsługuje kilka)
\end{frame}

\section{Rails}
\subsection{O co chodzi?}
\begin{frame}[fragile]
\frametitle{Zasady}
CoC, DRY, MVC
\end{frame}

\begin{frame}[fragile]
\frametitle{CoC}
\end{frame}

\begin{frame}[fragile]
\frametitle{MVC}
\end{frame}

\begin{frame}[fragile]
\frametitle{TDD?}
\end{frame}

\begin{frame}[fragile]
\frametitle{Jak zacząć?}

Heroku 
EngineYard
\end{frame}



\section{Podsumowanie}

\begin{frame}[fragile]
\frametitle{PRUG}
\end{frame}

\begin{frame}[fragile]
\frametitle{Big sites made with RoR}
Hulu, YellowPages, GitHub, Twitter (początkowo), Ask.fm
\end{frame}

\begin{frame}[fragile]
\frametitle{Rekrutujemy!}
\url{https://netguru.co/career}
\end{frame}
\end{document}